%% kapa: exam=1 round=35 year=2024 period=2 form=A subject=accounting subject_ko=회계학 seq=23 number=63
%% source: Q-Net 감정평가사 기출문제 (cst00309) artlSeq=5213229 file=raw/1008/5213229/2교시 문제지 원본.pdf
%% conversion: kordoc 4.24.51 -> tools/gen1.py (auto); review=proofread-sonnet; flags=tab
\begin{question}{63}
도소매업을 영위하는 (주)감평은 20×1년 초 건물을 취득(취득원가 ￦10,000, 내용연수 5년, 잔존가치 ￦0, 정액법 상각)하였다. 공정가치가 다음과 같을 때, (주)감평이 동 건물을 유형자산으로 분류하고 재평가모형을 적용하였을 경우(A)와 투자부동산으로 분류하고 공정가치모형을 적용한 경우(B), 20×2년 당기순이익에 미치는 영향은?
\par{\small\centering\begin{tabular}{|c|c|c|}\hline
구분 & 20×1년 말 & 20×2년 말 \tabularnewline\hline
공정가치 & ￦9,000 & ￦11,000 \tabularnewline\hline
\end{tabular}\par}
\smallskip
\noindent\begin{tabular}{@{}l@{\hspace{0.4em}}p{0.17\linewidth}p{0.17\linewidth}@{\hspace{1em}}l@{\hspace{0.4em}}p{0.17\linewidth}p{0.17\linewidth}@{}}
 & \uline{A} & \uline{B} & & \uline{A} & \uline{B} \tabularnewline
① & 영향없음 & ￦1,000 증가 & ② & ￦2,250 감소 & ￦1,000 증가 \tabularnewline
③ & ￦2,250 감소 & ￦2,000 증가 & ④ & ￦2,000 감소 & ￦2,000 증가 \tabularnewline
⑤ & ￦2,000 증가 & 영향없음 & & & \tabularnewline
\end{tabular}
\end{question}
\answer{3}
